\renewcommand{\theposttitle}{How fixing what's broke is not about knowing what to do}
\renewcommand{\thepostdate}{2024-01-02}
\vspace*{\fill}
\begin{center}
\begin{minipage}{0.92\textwidth}
\small
\textbf{Clear, Not Clever}

When designing a visual metaphor, one of the most important elements to nail down is the labels you use.

Labels are words, arrows, even punctuation that communicate what the metaphor is.

And the key to using effective labels is clarity. Clarity is step 1 of \textbf{my 3-Step framework to making a visual}.
\end{minipage}
\end{center}
\vspace{4mm}
\begin{center}
\includegraphics[width=\textwidth, height=0.55\textheight, keepaspectratio]{C:/Users/James/Documents/GitHub/VisualIdeasBook/data/images/c98fa1468cc374dab1f9b70f931c1bd6.png}
\end{center}
\vspace{2mm}
\begin{center}
\begin{minipage}{0.92\textwidth}
\small
What’s critical to remember—clarity is about removing as much friction as possible between you and your audience to get your message across.

(And that means you can’t hide behind fancy words if your message is unclear to begin with)

When you write in a fancy way to impress people, you make them do extra work just so you can seem cool

This issue of Visual I.D.E.A.s is sponsored by Superside:
\end{minipage}
\end{center}
\vspace*{\fill}